% Generated draft: requires mathematical and source review.
\documentclass[11pt]{article}
\usepackage[margin=25mm]{geometry}
\usepackage{fontspec}
\setmainfont{Latin Modern Roman}
\usepackage{amsmath,amssymb,mathtools,unicode-math}
\setmathfont{Latin Modern Math}
\usepackage{enumitem,needspace,xurl,hyperref}
\setlength{\parindent}{0pt}
\setlength{\parskip}{5pt}
\setlength{\emergencystretch}{3em}
\newcommand{\N}{\mathbb{N}}
\newcommand{\Z}{\mathbb{Z}}
\newcommand{\Q}{\mathbb{Q}}
\newcommand{\R}{\mathbb{R}}
\newcommand{\C}{\mathbb{C}}
\newcommand{\F}{\mathbb{F}}
\newcommand{\A}{\mathbb{A}}
\newcommand{\PP}{\mathbb{P}}
\newcommand{\E}{\mathbb{E}}
\newcommand{\eps}{\varepsilon}
\newcommand{\dd}{\,\mathrm{d}}
\begin{document}
\section*{UnsolvedMath ID: 20003122}
\textbf{Source problem code:} catalogue entry
\subsection{Definitions and mathematical statement}
D. Dumas asks how the least-Lipschitz function behaves for metrics beyond the hyperbolic locus, such as negatively curved Riemannian metrics or their closure. The displayed source omits a homotopy restriction on maps. Without one, constant maps make the infimum zero. In the marked interpretation used by the aid, for metrics on a fixed closed oriented manifold $M$ define
\[
L(g,h)=\inf_{\varphi\simeq\operatorname{id}_M}\operatorname{Lip}(\varphi:(M,g)\to(M,h)).
\]
The original Teichm\"uller setting concerns surfaces; some formulas below make sense in dimension $d$ more generally. Maps need not be diffeomorphisms for the infimum as written, and the marking restriction is explicit.

For positive-definite metrics on compact $M$, the aid asserts attainment of the infimum and
\[
L(g,k)\leq L(g,h)L(h,k),\qquad
L(g,h)\geq\sup_\alpha\frac{\ell_h(\alpha)}{\ell_g(\alpha)},\qquad
L(g,h)^d\operatorname{Vol}_g(M)\geq\operatorname{Vol}_h(M),
\]
where $\alpha$ ranges over the relevant nontrivial free homotopy classes. The volume inequality uses degree one. The summary reports existing work showing that equality with marked-length stretch, familiar for hyperbolic metrics, can fail in variable negative curvature; no example or full citation is supplied.

The proposed normalization is
\[
\widehat L(g,h)=L(g,h)\left(\frac{\operatorname{Vol}_g(M)}{\operatorname{Vol}_h(M)}\right)^{1/d}.
\]
It is at least $1$, submultiplicative, and unchanged by independently scaling $g$ and $h$: $L(a^2g,b^2h)=(b/a)L(g,h)$ while the volume factor cancels $b/a$. Thus $\widehat d(g,h)=\log\widehat L(g,h)$ is a nonnegative asymmetric metric kernel. The claimed zero-distance rigidity identifies metrics by a homothety and an isometry in the prescribed identity homotopy class. For closed negatively curved surfaces, this gives the reported asymmetric metric modulo $\operatorname{Diff}_0$ and homothety. In higher dimensions, identifying identity-homotopic diffeomorphisms with $\operatorname{Diff}_0$ needs an additional isotopy hypothesis and is not silently inferred.

For a uniform tensor-comparison distance, define
\[
\rho(g,g')=\inf\{r\geq0:e^{-2r}g\leq g'\leq e^{2r}g\}.
\]
The aid's quantitative stability statements are
\[
|\log L(g',h')-\log L(g,h)|\leq\rho(g,g')+\rho(h,h'),
\]
\[
|\log\widehat L(g',h')-\log\widehat L(g,h)|
\leq2\bigl(\rho(g,g')+\rho(h,h')\bigr).
\]
On a fixed-volume slice, the normalized logarithm equals the unnormalized one, so the first constant applies. These bounds provide continuous extension of the functionals to positive-definite $C^0$ metrics in the corresponding uniform topology. They do not assert an extension across degenerate or collapsed metrics, and metric separation on a nonsmooth closure is not concluded merely from continuity. The question's ``closure'' otherwise remains unspecified.
\subsection{Short English statement}
Interpret the omitted maps as marking-preserving ones; otherwise the infimum is zero. The supplied aid gives least-Lipschitz bounds, volume normalization and quantitative stability beyond hyperbolic metrics, while leaving degenerate closures and exact marked-length equality outside its scope.
\subsection{Sources}
\begin{itemize}
\item Local UnsolvedMath registry, permanent ID $20003122$, source code \texttt{catalogue entry}; canonical \texttt{aim-topology-notes.json}, \texttt{notes[209]}.
\item AIM \emph{Problems on Thurston Metric}, Problem 5.4, attributed to D. Dumas, workshop \emph{Lipschitz metric on Teichmueller space}: \url{https://aimath.org/pastworkshops/lipschitzteichproblems.pdf}. The supplied verification says the displayed PDF infimum contains no map restriction; the identity-homotopy condition is a conservative marked-space reconstruction. Statement status is \texttt{reconstructed\_unverified}; original remarks and literature field are empty.
\item The aid mentions existing variable-negative-curvature failures of marked-length equality without authoritative complete citations or explicit examples.
\item Research attempt 1: \texttt{partial\_result}, status \texttt{partially\_solved}, classification \texttt{PARTIAL-PROGRESS}, rigor \texttt{proved}, novelty confidence low. Its normalization, rigidity and stability assertions are machine-generated aid, not peer reviewed, and require independent verification before citation.
\end{itemize}
\end{document}
